% ==========================================================
% 110_conclusion.tex
% Forecast Admissibility Surface (FAS)
% ==========================================================

\section{Conclusion}
\label{sec:fas_conclusion}

Not all demand is forecastable, and not all slices of a forecast domain admit
stable learning, evaluation, or control. In operational systems, treating every
slice as equally admissible introduces noise, instability, and false confidence
that can undermine both analytic rigor and practical decision-making.

The Forecast Admissibility Surface addresses this gap by providing a
deterministic, slice-keyed governance construct that identifies where forecasting
is structurally justified and where it is not. By grounding admissibility in
baseline error anatomy, FAS exposes hostile regimes early in the forecasting
pipeline, before modeling complexity obscures structural signals.

Within the Forecast Readiness Framework, FAS complements existing diagnostics by
answering a foundational question that is often left implicit: \emph{Should this
slice be forecasted at all?} By making this decision explicit, auditable, and
joinable, FAS strengthens downstream modeling, evaluation, and governance without
introducing new optimization targets or subjective tuning.

More broadly, FAS reinforces a core principle of production forecasting: analytic
discipline begins not with better models, but with clear boundaries on where
models are appropriate. By enforcing those boundaries, the Forecast Admissibility
Surface contributes to more stable, interpretable, and trustworthy forecasting
systems.
